\medskip

On étudie l'évolution de la population d'une espèce animale au sein d'une réserve naturelle.


Les effectifs de cette population ont été recensés à différentes années. Les données collectées sont présentées dans le tableau suivant :

\begin{center}
	\begin{tabularx}{\linewidth}{|l|*{4}{>{\centering \arraybackslash}X|}} \hline
		Année              & 2000 & 2005 & 2010 & 2015 \\ \hline
		Nombre d'individus & 50   & 64   & 80   & 100  \\	\hline
	\end{tabularx}
\end{center}

Pour anticiper l'évolution de cette population, la direction de la réserve a choisi de modéliser le nombre d'individus en fonction du temps.

Pour cela, elle utilise une fonction, définie sur l'intervalle $[0 ~;~ +\infty [$, dont la variable $x$ représente le temps écoulé, en année, à partir de l'année 2000.

Dans son modèle, l'image de 0 par cette fonction vaut 50, ce qui correspond au nombre d'individus en l'an 2000.

	\subsection*{Partie A. Modèle 1}

	Dans cette partie, la direction de la réserve fait l'hypothèse que la fonction cherchée satisfait l'équation différentielle suivante :
	$$
	y'=0,05 y-0,5 \quad\left(E_{1}\right)
	$$

	\begin{enumerate}
		\item Résoudre l'équation différentielle $\left(E_{1}\right)$ avec la condition initiale $y(0)=50$.

		\item Comparer les résultats du tableau avec ceux que l'on obtiendrait avec ce modèle.
	\end{enumerate}

	\subsection*{Partie B. Modèle 2}

	Dans cette partie, la direction de la réserve fait l'hypothèse que la fonction cherchée satisfait l'équation différentielle suivante :
	$$
	y'=0,05 y(1-0,00125 y)
	$$

	On note $f$ la fonction définie sur $[0 ~;~ +\infty[$ par :
	$$
	f(x)=\dfrac{800}{1+15 \e^{-0,05 x}}
	$$

	et $C$ sa courbe représentative dans un repère orthonormé.

	\medskip

	À l'aide d'un logiciel de calcul formel, on a obtenu les résultats suivants.

	\textbf{Pour toute la suite de l'exercice, on pourra utiliser ces résultats sans les démontrer, sauf pour la question 5.}

	\begin{center}
	\renewcommand{\arraystretch}{2.5}
\begin{tabularx}{\linewidth}{|l|l|>{\centering \arraybackslash}X|} 
\hline
  & Instruction & Résultat \\	
\hline
1 & $f(x):=\dfrac{800}{1+15 \e^{-0,05 x}}$	& $f(x)=\dfrac{800}{1+15 \e^{-0,05 x}}$  \\	
\hline
2 & $f'(x):=$ Dérivée $(f(x))$ & $f'(x)=\dfrac{600 \e^{-0,05 x}}{\left(1+15 \e^{-0,05 x}\right)^{2}}$ \\ 
\hline
3 & $f''(x):=$ Dérivée ( $f'(x)$ ) & $f''(x)=\dfrac{30 \e^{-0,05 x}}{\left(1+15 \e^{-0,05 x}\right)^{3}}\left(15 \e^{-0,05 x}-1\right)$ \\ 
\hline
4 & Résoudre ( $15 \e^{-0,05 x}-1 \geqslant 0$ )  & $x \leqslant 20 \ln (15)$  \\		
\hline
\end{tabularx}
\renewcommand{\arraystretch}{1}













	\end{center}

	\begin{enumerate}
		\item Démontrer que la fonction $f$ vérifie $f(0)=50$ et que pour tout $x \in \R$ :
		$$
		f'(x)=0,05 f(x)(1-0,00125 f(x))
		$$

		On admet que cette fonction $f$ est l'unique solution de $(E_{2})$ prenant la valeur initiale de 50 en 0.

		\item Avec ce nouveau modèle $f$, estimer l'effectif de cette population en 2050. Arrondir le résultat à l'unité.

		\item Calculer la limite de $f$ en $+\infty$. Que peut-on en déduire quant à la courbe $C$ ? Interpréter cette limite dans le cadre de ce problème concret.

		\item Justifier que la fonction $f$ est croissante sur $[0 ~;~+\infty[$.

		\item Démontrer le résultat obtenu en ligne 4 du logiciel.

		\item On admet que la vitesse de croissance de la population de cette espèce, exprimée en nombre d'individus par an, est modélisée par la fonction $f'$.
		\begin{enumerate}
			\item Étudier la convexité de la fonction $f$ sur l'intervalle $[0 ~;~ +\infty[$ et déterminer les coordonnées des éventuels points d'inflexion de la courbe $C$.

			\item La direction de la réserve affirme :

			\og Au vu de ce modèle, la vitesse de croissance de la population de cette espèce va augmenter pendant un peu plus de cinquante ans, puis va diminuer \fg{}. La direction a-t-elle raison? Justifier.
		\end{enumerate}
	\end{enumerate}

\bigskip
